\section*{Chapter \ref{ch:g5:irreversible-changes} --- Changes That Cannot Be Undone}

\begin{solution}{exo:g5:irreversible-changes:1}
Melting chocolate: can be undone (it sets again when cooled). \omterm{def:g4:what-a-fire-needs:burning}{Burning} a
match: cannot. Dissolving salt: can be undone (evaporate the water).
Toasting bread: cannot.
\end{solution}

\begin{solution}{exo:g5:irreversible-changes:2}
A change in which new substances appear. In the kitchen: cooking an egg,
baking a cake (also: toasting bread, a cut apple browning).
\end{solution}

\begin{solution}{exo:g5:irreversible-changes:3}
By \omterm{def:g3:separating-mixtures:evaporate}{evaporating} the water: the salt stays in the dish. Dissolving is a
\omterm{def:g5:irreversible-changes:reversible}{reversible change}.
\end{solution}

\begin{solution}{exo:g5:irreversible-changes:4}
A new smell, and lumps (a solid) appearing in the milk.
\end{solution}

\begin{solution}{exo:g5:irreversible-changes:5}
The changes on the left: sugar dissolving in water, ice melting, sand
and gravel mixed. The double arrow means that the change can go both
ways: it can be undone.
\end{solution}

\begin{solution}{exo:g5:irreversible-changes:6}
The oil keeps \omterm{def:g4:air-a-mixture-of-gases:air}{air} and water away from the iron of the chain: it slows
down rusting.
\end{solution}

\begin{solution}{exo:g5:irreversible-changes:7}
Gas bubbles appearing in a liquid without any heating.
\end{solution}

\begin{solution}{exo:g5:irreversible-changes:8}
Useful: cooking food, baking bread, \omterm{def:g4:what-a-fire-needs:burning}{burning} gas for heating. Harmful:
rusting, food going off, wood rotting.
\end{solution}

\begin{solution}{exo:g5:irreversible-changes:9}
The cold slows down the \omterm{def:g5:irreversible-changes:chemical-change}{chemical changes} that make milk turn sour, so
it keeps longer.
\end{solution}

\begin{solution}{exo:g5:irreversible-changes:10}
No. The bubbles are water vapour, which turns back into water when it
cools: nothing new is made, and the change can be undone. Bubbles alone
do not prove that a new substance has appeared.
\end{solution}

\begin{solution}{exo:g5:irreversible-changes:11}
A \omterm{def:g4:what-a-fire-needs:burning}{burning} candle gives out heat and light, and makes new substances:
water vapour (a cold glass mists up above it) and carbon dioxide. The
wax that burns does not come back. Melting wax, on the other hand, sets
again when it cools.
\end{solution}
